%% ============================================================================
%%  Journal paper — Cross-Modal Generative Alignment with Explanation-Aware VLMs
%%
%%  Template : Elsevier `elsarticle` (Information Fusion / ESWA / Pattern Recog.)
%%  Status   : SKELETON. Methodology is drafted (stable, mirrors src/); Results
%%             and Discussion are placeholders filled as experiments land.
%%  Workflow : user runs experiments → hands over results.json → assistant
%%             journals them here. See README.md.
%%
%%  \note{...} = red reviewer note,  \todo{...} = blue action. Remove both before
%%  submission (comment out the two \newcommand lines to hide them).
%% ============================================================================
\documentclass[preprint,review,12pt]{elsarticle}

\usepackage[T1]{fontenc}
\usepackage[utf8]{inputenc}
\usepackage{amsmath,amssymb}
\usepackage{graphicx}
\usepackage{booktabs}
\usepackage{multirow}
\usepackage{xcolor}
\usepackage{url}
\usepackage[hidelinks]{hyperref}

\newcommand{\note}[1]{\textcolor{red}{[\textbf{NOTE:} #1]}}
\newcommand{\todo}[1]{\textcolor{blue}{[\textbf{TODO:} #1]}}

\journal{Information Fusion} % alt: Expert Systems with Applications

\begin{document}

\begin{frontmatter}

\title{Cross-Modal Generative Alignment for Document, Chart, and Image
Reasoning with Explanation-Aware Vision--Language Models}

\author[a]{Komron Valijonov}
\author[a]{Raja Hashim Ali}
\affiliation[a]{organization={\note{Department / University}},
  city={\note{city}}, country={\note{country}}}

\begin{abstract}
\note{Rewrite with real numbers once results land. Draft below adapts the proposal abstract.}
Vision--language models (VLMs) achieve cross-modal reasoning by aligning visual
features with pretrained language models, yet two design choices remain
underexplored for \emph{structured} visual content such as documents and charts:
(i) the \emph{alignment objective} --- contrastive versus generative --- and
(ii) the effect of \emph{explanation-aware} training on not only answer accuracy
but the \emph{faithfulness} of generated rationales. We present a controlled
study that fixes backbone, data, and budget and varies only the training signal,
across four alignment variants (baseline, contrastive-enhanced, generative, and
explanation-aware generative) fine-tuned with QLoRA on a single 40\,GB GPU. We
evaluate on DocVQA, ChartQA, ScienceQA, A-OKVQA, and VQAv2 along answer accuracy,
explanation quality (BLEU, ROUGE-L, BERTScore), and faithfulness
(attention-alignment and evidence-masking consistency). \todo{headline findings
+ numbers}. Our results indicate \todo{one-sentence claim: explanation-aware
QLoRA SFT improves faithfulness, and the gain is domain-dependent}.
\end{abstract}

\begin{keyword}
vision--language models \sep cross-modal alignment \sep explanation-aware training
\sep chart and document understanding \sep faithfulness \sep QLoRA
\end{keyword}

\end{frontmatter}

% ============================================================================
\section{Introduction}
\label{sec:intro}
\note{Motivation $\rightarrow$ problem $\rightarrow$ contributions $\rightarrow$ RQs $\rightarrow$ outline.}
\todo{Sharpen the headline claim (see README): cheap explanation-aware QLoRA
\emph{supervised} fine-tuning improves rationale \emph{faithfulness}, not just
accuracy, and the benefit is larger for structured visual content (charts,
documents) than for natural images --- a gap left open by RL-based and
evaluation-only prior work.}

\paragraph{Contributions.}
\begin{itemize}
  \item \todo{C1 --- controlled comparison of contrastive vs.\ generative vs.\ explanation-aware alignment under a fixed single-GPU budget (RQ2, RQ3).}
  \item \todo{C2 --- faithfulness analysis of generated rationales for structured visual reasoning (RQ5).}
  \item \todo{C3 --- domain-generality / transfer profile of explanation-aware alignment (RQ4, RQ7).}
  \item \todo{C4 --- a reproducible, registry-driven QLoRA recipe (released, Apache 2.0).}
\end{itemize}

% ============================================================================
\section{Related Work}
\label{sec:related}
\note{RQ1 survey lives here. Drafted in parallel with experiments.}
\todo{Threads to cover and the gap to claim:}
\begin{itemize}
  \item \emph{Alignment paradigms} --- contrastive (CLIP~\cite{clip}) vs.\ generative (BLIP-2~\cite{blip2}); instruction-tuned VLMs (Qwen2-VL~\cite{qwen2vl}, PaliGemma~\cite{paligemma}).
  \item \emph{Explanation / chain-of-thought training} --- ScienceQA~\cite{scienceqa}, Multimodal-CoT~\cite{multimodalcot}.
  \item \emph{Faithfulness \& hallucination} --- the NLE-faithfulness benchmark~\cite{nlefaithfulness}, counterfactual testing~\cite{edct}; object hallucination (POPE~\cite{pope}).
  \item \emph{Structured-visual faithfulness} (closest neighbors) --- ChartHal~\cite{charthal} (evaluation only), Chart-RVR~\cite{chartrvr} (RL-based). \textbf{Gap:} no controlled study of explanation-aware \emph{SFT} on faithfulness \emph{across} document/chart/natural-image domains.
  \item \emph{Parameter-efficient fine-tuning} --- LoRA~\cite{lora}, QLoRA~\cite{qlora}.
\end{itemize}

% ============================================================================
\section{Methodology}
\label{sec:method}
% ---- DRAFTED: stable, mirrors the released src/ framework. ----

\subsection{Data sources and preprocessing}
\label{sec:data}
We use five public benchmarks spanning three visual domains: documents
(DocVQA~\cite{docvqa}), charts (ChartQA~\cite{chartqa}), and natural images
(ScienceQA~\cite{scienceqa}, A-OKVQA~\cite{aokvqa}, VQAv2~\cite{vqav2}).
Explanation-aware supervision requires gold rationales: ScienceQA and A-OKVQA
provide them natively; for the remaining datasets rationales are
\todo{synthesized via a teacher model and filtered, or treated as answer-only}.
All samples are mapped to a single neutral schema (image, question, answer,
optional choices, optional rationale). \note{Report the filtering yield, e.g.\
of 12{,}726 ScienceQA training items $\approx$5k retain an image + gold rationale.}

\subsection{Model design: four alignment variants}
\label{sec:variants}
Holding the backbone fixed (Qwen2-VL-2B-Instruct~\cite{qwen2vl} primary; BLIP-2~\cite{blip2}
and PaliGemma~\cite{paligemma} for cross-architecture breadth), we compare four
training signals that differ \emph{only} in the objective:
\begin{enumerate}
  \item \textbf{Baseline} --- native projection, no extra alignment loss.
  \item \textbf{Contrastive-enhanced} --- baseline + auxiliary InfoNCE between
        projected visual features and answer sentence embeddings.
  \item \textbf{Generative} --- next-token prediction on $(\text{image}, Q \rightarrow A)$.
  \item \textbf{Explanation-aware generative} --- next-token prediction on
        $(\text{image}, Q \rightarrow \text{explanation} \rightarrow A)$.
\end{enumerate}

\subsection{Explanation-aware objective}
\label{sec:loss}
The explanation-aware loss weights the answer and explanation spans of the target:
\begin{equation}
  \mathcal{L} \;=\; \alpha \, \mathcal{L}_{\text{answer}}
              \;+\; (1-\alpha)\, \mathcal{L}_{\text{explanation}},
  \label{eq:loss}
\end{equation}
where each term is the masked cross-entropy over its token span and
$\alpha \in \{0.0, 0.5, 1.0\}$ ($\alpha{=}1$ recovers answer-only / generative;
$\alpha{=}0$ is explanation-only). Spans are tagged at the token level so the two
losses are computed from a single forward pass.

\subsection{Parameter-efficient fine-tuning (QLoRA)}
\label{sec:qlora}
All variants are fine-tuned with QLoRA~\cite{qlora}: the base model is quantized
to 4-bit NF4 and frozen, and low-rank adapters (rank 16, $\alpha{=}32$) are
trained on the language-model attention/MLP projections and the alignment module;
the vision encoder is frozen. Optimizer AdamW ($\beta{=}(0.9,0.999)$, weight
decay $0.01$), learning rate $2\times10^{-4}$ with cosine schedule and 3\%
warm-up, effective batch size 32 via gradient accumulation, bf16. \note{1 seed
per configuration --- all claims carry a variance caveat.}

\subsection{Evaluation metrics}
\label{sec:metrics}
\textbf{Answer accuracy:} exact match / ScienceQA \& A-OKVQA accuracy, DocVQA
ANLS, ChartQA relaxed accuracy. \textbf{Explanation quality:} BLEU-4, ROUGE-L,
BERTScore-F1 (generated rationale vs.\ gold). \textbf{Faithfulness (RQ5):}
attention-alignment between image-patch attention and explanation-referenced
regions, and evidence-masking consistency (answer-likelihood drift when salient
regions are occluded). Multiple-choice answers are scored by a robust
generate-then-score (chain-of-thought) protocol rather than string matching.

% ============================================================================
\section{Experimental Setup}
\label{sec:setup}
\note{Backbones, dataset splits/sizes, the experiment grid, compute (1$\times$A100-40GB),
software (PyTorch, Transformers, PEFT, BitsAndBytes), and the released code.}

% ============================================================================
\section{Results}
\label{sec:results}
\note{Filled from results.json as runs land. One subsection per RQ; tables below
are empty scaffolds with their captions already fixed.}

\subsection{RQ2 --- Contrastive vs.\ generative alignment}
\begin{table}[t]
  \centering
  \caption{Answer accuracy by alignment strategy across domains.}
  \label{tab:rq2}
  \begin{tabular}{lccccc}
    \toprule
    Variant & DocVQA & ChartQA & ScienceQA & A-OKVQA & VQAv2 \\
    \midrule
    Contrastive-enhanced & -- & -- & -- & -- & -- \\
    Generative           & -- & -- & -- & -- & -- \\
    \bottomrule
  \end{tabular}
\end{table}

\subsection{RQ3 --- Explanation-aware training}
\note{Accuracy + explanation-quality table; the $\alpha$-sweep curve
($\{0,0.5,1\}$). Lead with Qwen2-VL (BLIP-2's BLEU/ROUGE are near-zero --- weak generator).}

\subsection{RQ4 --- Architectural / training factors}
\subsection{RQ5 --- Faithfulness and hallucination}
\subsection{RQ6 --- Parameter-efficient vs.\ scale}
\subsection{RQ7 --- Cross-domain transfer}

% ============================================================================
\section{Discussion}
\label{sec:discussion}
\note{Interpret across RQs; the domain-generality hypothesis; limitations
(1 seed, QLoRA-only, 2B/7B scale, synthesized rationales).}

% ============================================================================
\section{Conclusion}
\label{sec:conclusion}
\note{Answer each RQ in one line; restate the contribution; future work
(more seeds, full fine-tuning, larger backbones).}

% ============================================================================
\section*{Acknowledgments}
\note{Supervisor, compute, funding if any.}

\bibliographystyle{elsarticle-num}
\bibliography{references}

\end{document}
